\subsection{Spherical Fourier-Bessel Basis Functions}\label{subsec:sfb_basis}

This section introduces the complete orthonormal basis of spherical Fourier-Bessel functions and their key properties.

\subsubsection{The SFB basis set}

The SFB basis functions are products of radial and angular components:

\begin{equation}
\psi_{\ell m}^{(n)}(\mathbf{r}) = A_{\ell n} j_\ell(k_{\ell n} r) Y_{\ell m}(\theta, \phi)
\label{eq:sfb_basis_function}
\end{equation}

where:
- $j_\ell$ is the spherical Bessel function of order $\ell$
- $Y_{\ell m}$ is the spherical harmonic
- $k_{\ell n} = x_{\ell,n} / R$ with $x_{\ell,n}$ the $n$-th zero of $j_\ell$
- $R$ is the domain radius (often taken to infinity)
- $A_{\ell n}$ is a normalization constant

\subsubsection{Domain and boundary conditions}

For cosmological fields defined on $r \in [0, \infty)$:
- Regularity at origin: $j_\ell(0) = 0$ for $\ell > 0$ ✓
- Decay at infinity: $j_\ell(kr) \to \cos(kr - \ell\pi/2)/r$ ✓

This makes Bessel functions natural for fields with no sources at infinity.

\subsubsection{Advantages over plane waves}

\begin{itemize}
    \item **Spherical symmetry**: Respects the observer-centric geometry of surveys
    \item **Finite domains**: Natural handling of radial selection functions
    \item **Boundary conditions**: Automatic regularity at origin and decay at large $r$
    \item **Angular locality**: Spherical harmonics separate angular and radial information
\end{itemize}

See \cref{subsec:sfb_orthonormality} for orthonormality relations and \nameref{subsec:sfb_expansion} for field expansions.
